\section{S9: Light's Requirement on the Medium --- Shear on the Brane, None in
the Bulk}\label{step:S9}

\stagefield{Part / anchor}{Part I --- Light. Sector~1, step~1.}
\claimstatus{\StatusOpen}
\stagefield{Purpose}{State what light needs the medium to have. Light is a
transverse wave and a GNLS superfluid cannot carry one, so light's first act is
to name the structure the GNLS does not supply. The step introduces the two
constants the whole sector rests on and fixes the light cone as their ratio.}

\paragraph{Status.}
\StatusOpen. First step banked under the requirements-first direction: light
states what it needs, and the medium's structure is not assumed in advance. The
result is a \emph{requirement with a live falsifier}, not a derivation --- both
constants enter \textbf{postulated}, and whether one substructure can meet both
halves of the requirement is a knit question this step does not settle.

\paragraph{The argument, in three moves.}
\resultanchor{1 --- What Maxwell demands.} Two transverse polarisations, one
speed, no longitudinal mode.

\resultanchor{2 --- What a GNLS has.} Linearising Madelung about a uniform
\(\rho_0\) gives exactly \emph{one} mode,
\[
  \omega^2 = c_s^2k^2 + \frac{\hbar^2}{4m^2}k^4,
  \qquad
  c_s^2 = \frac{1}{m}\frac{dP}{d\rho} = \frac{5K\rho_0^4}{m},
\]
and it is \textbf{longitudinal}. The obstruction is not incidental but
representation-theoretic: \(\Psi=\sqrt{\rho}\,e^{i\theta}\) is one complex
scalar, its excitations \(\delta\rho\) and \(\delta\theta\) are both spin-0, and
one broken \(U(1)\) gives one Goldstone phonon --- so a scalar's spectrum cannot
contain a spin-1 photon, \emph{in any dimension}.

\resultanchor{3 --- Therefore.} Light requires structure the GNLS does not
contain. Brane elasticity is not a consequence of the GP/NLS mean field: a
GP/NLS superfluid is a \emph{fluid}, with zero shear modulus. The brane's shear
rigidity therefore requires the \textbf{substructure} beneath the mean field ---
the standing postulate that the GNLS is the effective, coarse-grained
description of a medium with constituents of its own.

\paragraph{Provenance of moves 2 and 3 --- one source, not two.}
Both quotations come from the same document,
\StageFile{software/stage1\_solver/decisions/15\_em\_medium\_native\_physical\_picture.md}
(lines 29--40 and 230--233). It is an external review, and \textbf{no script in
this repository executes either argument}. A weaker in-repo form exists ---
\(v=(\hbar/m)\nabla\theta \Rightarrow \nabla\times v\equiv 0\), so superfluid
flow is curl-free --- but it is never wired into a no-shear proof. The claim
that a GNLS cannot carry light therefore rests on a single external document
with no executing script; it is stated here at that strength and no higher.

\paragraph{The requirement --- and it is two-sided.}
\begin{center}
\begin{tabular}{@{}lll@{}}
\toprule
phase & requirement & what it buys \\
\midrule
ordered (\(\chi_B=1\))    & carries MacCullagh curl-only shear & photons \emph{exist} \\
disordered (\(\chi_B=0\)) & carries \emph{no} shear            & photons stay \emph{confined} \\
\bottomrule
\end{tabular}
\end{center}

\noindent
Both halves are light's own bill. The second is \emph{not} magnetism's: without
it a photon has somewhere else to go, light leaks off our three-dimensional
space, and that is an energy sink with no observational room. The requirement is
stressed hardest at the throat, where the brane is bent into \(\pm w\).

A \emph{third} consumer of the second half: the recorded throat-support
mechanism is an outward trapped standing-wave pressure whose trapped mode is a
brane-shear standing wave. If the bulk carried shear, that mode radiates away,
the outward pressure vanishes, and the throat closes. Bulk shear-freeness is
what makes the geon possible at all.

\paragraph{The equations.}
\[
  L=\tfrac12\rho_{br}(\partial_t u)^2-\tfrac12\mu_R(\nabla\times u)^2,
  \qquad
  c_\gamma^2=\frac{\mu_R}{\rho_{br}} \quad (R4).
\]

\paragraph{Dimensions, derived from the action.}
\(L\) is an energy density on the three-dimensional brane,
\([L]=M\,L^{-1}T^{-2}\), and \([u]=L\); hence \([\partial_t u]=LT^{-1}\) and
\([\nabla\times u]=1\), so
\[
  [\rho_{br}]=M\,L^{-3}=[-3,0,1],
  \qquad
  [\mu_R]=M\,L^{-1}T^{-2}=[-1,-2,1],
\]
and \(R4\) gives \([-1{-}({-}3),\,-2{-}0,\,1{-}1]=[2,-2,0]=2\times[1,-1,0]\), so
\(c_\gamma=[1,-1,0]\). These were derived from the Lagrangian \emph{before}
consulting the parameter register, which records the same two dimensions
independently (\StageFile{notes/parameter\_register.md}, lines 137--138).

\paragraph{What is new.}
\begin{center}
\begin{tabular}{@{}lll@{}}
\toprule
item & class & why \\
\midrule
\(\rho_{br}\) & \textbf{postulated} & a substructure property; not derivable from the GNLS \\
\(\mu_R\)     & \textbf{postulated} & ditto; the polar-\(P\) route returned \texttt{FAIL\_COUPLE\_STRESS\_NOGO} (B2) \\
\(c_\gamma\)  & \textbf{derived}    & \(R4\), from the two above \\
curl-only form & \textbf{postulated} & forced by Maxwell's no-longitudinal demand \\
\bottomrule
\end{tabular}
\end{center}

\noindent
The far field consumes only the ratio \(c_\gamma\). That must not be read as
``the two are not separately owed'': a \emph{simulation} needs \(\mu_R\) and
\(\rho_{br}\) absolutely.

\paragraph{Inputs, by locus.}
The standard GNLS (\(\psi\), \(\rho\), \(\theta\), Madelung,
\(v=(\hbar/m)\nabla\theta\)) comes standard with the vocabulary; Maxwell's three
demands come from the imported theory. \textbf{Nothing is taken from an earlier
v3 step} --- S9 is sector~1, step~1.

\paragraph{Regime.}
Linear response about an unstrained brane; small oscillations. Everything here
is quadratic.

\paragraph{Departure.}
None at this step. The departure enters at S11, when the medium's
compressibility overrides the no-longitudinal demand.

\paragraph{Falsifier, and its status.}
\emph{Requirement:} the substructure carries shear in its ordered phase and none
in its disordered phase. \emph{Falsifier:} a substructure supplying one and not
the other --- in particular one carrying shear in the \emph{disordered} phase,
which would unconfine light \emph{and} kill the trapped-mode geon.
\emph{Status:} \textbf{live}, and it is a knit question. Three independent
consumers (photon propagation, photon confinement, geon support), and nothing
yet shows one substructure can do both. Bulk shear-freeness is currently
\textbf{postulated}; the only machine check on it is dimensional.

This is light's first open requirement conflict, and it is \emph{internal to
light} --- not a light-versus-magnetism conflict. It is a question about the
ordering transition: can one substructure be ordered-and-shear-bearing in one
phase and unstructured-and-shear-free in the other?

\paragraph{Registry additions executed at this step.}
\(Q.brane.rho\_br\ [-3,0,1]\), \(Q.brane.mu\_R\ [-1,-2,1]\),
\(Q.brane.c\_gamma\ [1,-1,0]\), and relation \(R4\):
\(c_\gamma=\sqrt{\mu_R/\rho_{br}}\). \(R4\) is the corpus's existing name for
this relation. Note that \(c_\gamma\) re-enters \emph{here}, with provenance,
having been deleted from the \emph{medium} block at S0.5: that round trip is the
mechanism working.

% NOT \stagefield{Verification} -- that field is suppressed in the default
% build, and this is a GAP DISCLOSURE, not a verification listing.  It must
% remain visible.  Delete these two artifacts from the "owed" wording only when
% they actually exist.
\paragraph{Verification.}
The cone was derived independently \textbf{four} times and agrees each time:
the registry entry (walked with the user), a \emph{blind} Mathematica audit
(\StageFile{research/pde\_ledger\_v3/mathematica/S9\_light\_requires\_shear\_mathematica\_audit.wl}),
and two independent reviewers deriving it by hand. All four give
\(c_\gamma^2=\mu_R/\rho_{br}\), \([\rho_{br}]=\{-3,0,1\}\) and
\([\mu_R]=\{-1,-2,1\}\).

The Mathematica audit is blind by design: it takes no arguments, imports
nothing, and does not read the registry --- so it shares no input with the
thing it is compared against, and the agreement is not vacuous. It derives
rather than asserts (varies the action, forms the dynamical matrix, solves
\(\det=0\)); no dimension vector and no speed formula appears as a literal in
its source. Registry gates: dimensional homogeneity \StatusText{pass},
acceptance \StatusText{match}, able-to-fail \StatusText{pass}.

\emph{Scope of that audit, stated so it is not over-read.} Its own
\texttt{PASS} means its internal consistency checks did not contradict one
another. It is insensitive to a wrong overall prefactor, a flipped potential
sign, and the assumed brane dimension --- each of which changes what it prints
while leaving its verdict green. Those are caught by the external comparison
above, which is therefore load-bearing rather than a formality.

\emph{No separate SymPy audit was written for this step, deliberately.} The
derivation is two lines of algebra already executed by the Mathematica engine
and by v2's script
(\StageFile{research/pde\_ledger\_v2/scripts/ledger\_stage003\_transverse\_photons\_stray\_longitudinal\_sympy\_audit.py},
lines 516--601), cited by locus. A second engine earns its place where the
algebra is long enough that it could genuinely disagree; here it would be a
third execution of the same two lines.
